\chapter{多层感知机与神经网络入门}
\chaptergoal{能够把逻辑回归看作没有隐藏层的概率分类模型，说明输入层、隐藏层和输出层的计算；理解激活函数、损失、反向传播与优化器的不同职责；使用 sklearn 的 MLPClassifier 和 MLPRegressor；通过对照实验分析标准化、隐藏层、正则化、训练轮数和随机初始化的影响。}

前面学习的线性模型、树、近邻和核方法各有表达能力与训练方式。多层感知机（multilayer perceptron, MLP）提供另一种途径：让模型通过多次加权求和和非线性变换，形成新的输入表示。它是最基本的前馈神经网络之一，也是从经典机器学习走向后续神经网络课程的自然桥梁。本章不是要把深度学习的所有细节讲完，而是要使学生能够真正理解一个小型网络的计算与训练。
\section{从逻辑回归到隐藏层}

\textbf{人工神经元}在本章只是一个计算单元：接收若干数值，各乘一个权重，加上偏置，再按规定函数产生一个输出。这个名称借用了神经元的启发，但这里不把它当成对真实生物神经细胞的完整模型。多个单元并列组成一层，前一层的输出作为后一层的输入。

\textbf{输入层}记录原始特征，\textbf{输出层}产生任务要求的预测；两者之间的计算层称为\textbf{隐藏层}。“隐藏”表示它不是直接给定的输入或监督目标，不是说程序无法查看它。\textbf{前馈}表示信号沿输入到输出的方向计算，不在这一次计算中回到前面的节点形成循环。\textbf{激活函数}就是加权求和之后施加的数值变换，稍后通过具体计算说明其作用。

第6章的二分类逻辑回归先计算 $z=\mathbf w^\mathsf T\mathbf x+b$，再通过 Sigmoid 输出概率。这已经包含输入、权重、偏置和非线性输出。如果只考虑输出节点，它可以表示为一个\textbf{没有隐藏层}的简单网络。这样的网络在原始特征空间中的决策边界仍由线性得分决定。

现在考虑先对输入做一次变换，再根据变换后的结果进行分类：
\begin{align}
\mathbf h&=\phi(W_1\mathbf x+\mathbf b_1),\\
z&=\mathbf w_2^\mathsf T\mathbf h+b_2,\qquad p=\sigma(z).
\end{align}
这里 $\phi$ 表示隐藏层激活函数，$\mathbf h$ 是隐藏层产生的新表示。若 $\mathbf x$ 有 $p$ 个特征、隐藏层有 $q$ 个单元，则 $W_1$ 的维度为 $q\times p$。理解这种矩阵尺寸有助于学生判断代码参数的含义，而不必先完成复杂的矩阵求导。

\fig{ch12_network.pdf}{逻辑回归的直接概率输出与单隐藏层 MLP 的结构对比。隐藏层使模型可以组合输入形成新的非线性表示。}

\subsection{表示能力：一个线性分类器确实做不到的任务}
考虑四条人工记录：$(0,0)$与$(1,1)$属于类别0，$(1,0)$与$(0,1)$属于类别1。这称为异或关系，表示两个输入不同则输出1。设线性得分为 $w_1x_1+w_2x_2+b$，类别1需正得分。若四点都严格分对，则 $b<0$、$w_1+b>0$、$w_2+b>0$、$w_1+w_2+b<0$。

中间两式相加得到 $w_1+w_2+2b>0$，因此 $w_1+w_2+b>-b>0$，与最后一式矛盾。这一小证明表明，无论线性模型迭代多少次，都不能用一条直线严格分开四点。问题在表示能力，不是学习率不够好。

\worked{两个隐藏单元怎样表示异或}
定义
\[
h_1=\max(0,x_1-x_2),\qquad h_2=\max(0,x_2-x_1),\qquad s=h_1+h_2.
\]
对 $(0,0),(1,0),(0,1),(1,1)$，隐藏表示分别为 $(0,0),(1,0),(0,1),(0,0)$，输出 $s$依次为0、1、1、0。使用阈值0.5即可完全分类。这个网络的参数由我们为解释机制而手工指定，不是训练结果；它证明这种结构\emph{能够}表示关系，不证明任意随机初始化与优化预算都会学到相同解。

这个区别会贯穿神经网络学习：\textbf{存在合适参数}、\textbf{优化器能找到它}与\textbf{新数据上有效}是三个不同问题。增加隐藏层改善了可表示函数集合，但优化与泛化仍需实验检验。

\subsection{隐藏表示为什么由最终任务驱动}
逻辑回归直接对原始特征加权，MLP先用若干有参数的变换构造中间特征，再把这些新特征交给输出层。训练数据通常没有给每个隐藏单元规定目标值；隐藏层通过它对最终损失的影响得到调整。这称为\textbf{表示学习}，不是隐藏单元天然具有“温度”“亮度”这样的物理意义。

在葡萄酒化学分类中，一个隐藏单元可能组合几项成分，也可能只是为最终边界形成有用的数值方向。可解释它的数值计算，但不能未经证据给每个单元命名为某种真实化学机制。后续更复杂网络也是如此，结构复杂不自动带来因果解释。

\section{权重、偏置与前向计算}
一层网络通常先进行仿射变换 $\mathbf z^{(\ell)}=W^{(\ell)}\mathbf a^{(\ell-1)}+\mathbf b^{(\ell)}$，再通过激活函数得到 $\mathbf a^{(\ell)}=\phi(\mathbf z^{(\ell)})$。输入层的激活就是原始输入 $\mathbf a^{(0)}=\mathbf x$。一层层从输入计算到输出，称为\textbf{前向传播}。

\subsection*{把一个小网络从头算到尾}
取人工输入$\mathbf x=(1,2)^{\mathsf T}$，两个隐藏单元的参数为
\[
W_1=\begin{pmatrix}1&0\\-1&1\end{pmatrix},\qquad
\mathbf b_1=\begin{pmatrix}0\\-2\end{pmatrix}.
\]
第一个单元的加权和为$1\times1+0\times2+0=1$；第二个为$-1\times1+1\times2-2=-1$。采用$\mathrm{ReLU}(z)=\max(0,z)$，即正数保留、负数变为零，两个输出便是$\mathbf h=(1,0)^{\mathsf T}$。

输出层权重为$(2,-1)$、偏置为$-0.5$，因此得分为$2\times1-1\times0-0.5=1.5$。用第6章的Sigmoid转成类别1的概率，得到$1/(1+e^{-1.5})\approx0.8176$。按0.5阈值预测为类别1。这些参数只是为手算指定的，尚未通过任何真实数据训练。

\fig{v2_mlp_forward.pdf}{两输入、两隐藏单元、一个输出的小网络数值计算。箭头表示计算顺序，框内给出当前中间数值。}

这个网络需要学习的量有：隐藏层4个权重和2个偏置，输出层2个权重和1个偏置，总计9个参数。输入维数2和隐藏单元数2不是这9个参数中的一部分；它们决定参数数组的形状。

权重和偏置属于模型参数；隐藏层单元数、隐藏层数量、激活函数选择等属于超参数。前向传播利用当前参数计算预测，并不自动说明参数怎样被学习。为了训练网络，还需要规定损失以及用于求梯度和更新参数的方法。

\subsection{先数清形状，再进行矩阵计算}
对一条含 $p$ 个特征的列向量 $\mathbf x$，含 $q$ 个隐藏单元的层可写为
\[
\mathbf z=W_1\mathbf x+\mathbf b_1,\quad \mathbf h=\phi(\mathbf z),
\]
其中 $W_1$ 为 $q\times p$，$\mathbf b_1$ 为 $q$ 维，激活逐元素计算。若输出有 $K$ 个得分，则 $W_2$ 为 $K\times q$，$\mathbf b_2$ 为 $K$ 维，输出得分为 $W_2\mathbf h+\mathbf b_2$。矩阵的行列取决于本书选择的列向量约定；程序常把一批样本按行存放，需要相应转置，不能仅凭矩阵乘法没有报错就认为含义正确。

\worked{一个葡萄酒分类网络有多少可学习量}
输入13维，隐藏层8个单元，输出3类得分。第一层有 $13\times8=104$ 个权重和8个偏置；输出层有 $8\times3=24$ 个权重和3个偏置，总数为139。改变训练样本数不会改变这个结构的参数个数，但改变隐藏宽度会。输出概率经Softmax由得分得到，不额外增加一组可学习参数。

当输入是一批 $n$ 条记录时，特征矩阵为 $n\times13$，隐藏输出为 $n\times8$，最后得分为 $n\times3$。同一组权重作用于每一行，这就是模型能够处理不同数量新样本的原因。它不为每条训练记录分别保存139个参数。

\subsection{偏置、中心化与数据流}
偏置允许激活阈值相对输入原点移动。例如 $h=\max(0,2x-1)$ 在 $x=0.5$处发生折点，而没有偏置的 $\max(0,2x)$ 折点固定在0。输入中心化会改变表示的原点，却不等于可以删除所有层的偏置；网络内部各层的分布仍会变化。

前向计算始终按固定顺序执行。中间输出既是下一层的输入，也是之后计算梯度所需的信息。只保存最后预测，可能不足以高效求出每层导数；反向传播会复用这些中间量。这是把程序按计算图组织的直接原因。

\section{激活函数与非线性表示}

例如第一层计算$h=2x+1$，第二层计算$\hat y=3h-4$。代入后得到$\hat y=6x-1$，仍是一条直线。层数增加了，能够表示的关系没有因这两层而变成弯曲关系。若第一层改为$h=\max(0,2x+1)$，则$x\leq-0.5$与$x>-0.5$时使用不同表达式，才引入了非线性的分段关系。这里“非线性”描述整体输入与输出关系，不意味着每一步都需要复杂函数。

如果每层只有线性变换，没有非线性激活，那么多个线性层的组合仍等价于一个线性变换；增加层数并不会赋予模型表达非线性边界的能力。常见的 ReLU 为 $\mathrm{ReLU}(z)=\max(0,z)$，Sigmoid 为 $1/(1+e^{-z})$。不同激活函数影响梯度传播、优化稳定性和表示能力。

ReLU 对正输入保留线性值，对负输入输出零；它计算简单，在许多神经网络中很常见，但也可能出现某些单元长期不激活的问题。Sigmoid 把结果压缩到0和1之间，适合二分类概率输出的解释，但在输入绝对值很大时导数较小，可能不利于深层网络中的梯度传播。激活函数的选择要与网络结构和任务共同理解。

\subsection{非线性对表示与训练的双重作用}
没有激活的两层仿射变换满足 $W_2(W_1\mathbf x+\mathbf b_1)+\mathbf b_2=(W_2W_1)\mathbf x+(W_2\mathbf b_1+\mathbf b_2)$，仍可合并为一层。加入ReLU后，不同输入可能激活不同的单元集合，在各区域中表现为不同的线性表达。这解释了一个有限ReLU网络为什么可以拼出较复杂的分段关系。

但激活函数不仅影响可表示函数，也影响梯度。Sigmoid导数为 $\sigma(z)[1-\sigma(z)]$，在0处为1/4，在绝对值很大的区域趋近0。若链式计算连续乘上许多很小的导数，靠前层得到的梯度可能很小。ReLU在正区间导数1、负区间导数0，减少了某些饱和问题，却可能使长期落在负区间的单元几乎不更新。不存在对所有任务都最好的激活。

\worked{观察一个单元是否还有学习信号}
设 $h=\max(0,wx+b)$，$x=2,w=-1,b=0$，则预激活为 $-2$，输出0，ReLU导数为0。无论后面损失对 $h$ 的导数当前是多少，通过该单样本回传到 $w$ 的梯度都乘以0。若所有训练输入都让它落在负区间，这个单元可能长时间不起作用。

这并不意味着任何一次负输入都是错误。不同样本可以激活不同单元，稀疏活动本来就是模型的一种表示方式。诊断应看整个训练过程与样本范围，而不是看到某个隐藏输出为0就改变激活函数。

\subsection{输入输出的物理约束不能靠激活名称保证}
回归输出层常使用线性输出，可以产生任意实数。若目标是非负量，必须检查模型是否出现负预测，考虑目标变换或符合任务的输出约束；简单裁剪会改变误差分布，也应记录。分类输出的概率和为1只是数学规范，不说明概率已经校准，更不说明输入样本一定属于训练类别之一。

\section{损失函数与网络训练}

\textbf{初始化}是开始训练前给权重和偏置设定起始值；\textbf{一个训练轮次（epoch）}通常表示把训练记录遍历一遍；\textbf{小批量（mini-batch）}是一次计算梯度时使用的一部分记录。若有100条记录、每批20条，一轮通常包含5次参数更新。因此“训练100轮”与“更新100次”未必相同。

网络中的\textbf{表示学习}是指隐藏层变换本身的参数也通过最终任务损失学习。葡萄酒分类数据只给最终类别，并没有给出“第3个隐藏单元应该等于几”的答案。训练通过损失与各层参数的联系调整这些中间表示。

网络训练必须明确任务。对于二分类，可通过概率输出与交叉熵度量差异；多分类常使用 Softmax 概率与交叉熵；回归可以使用平方误差等连续目标损失。正则化则通过限制过大的权重来控制复杂度。训练目标与独立评价指标可能不同：以交叉熵训练的分类器仍需用召回率、精确率或其他任务指标来评价。

训练的一次迭代大体包括：输入一批训练样本；前向计算预测；计算训练损失；求损失对权重与偏置的梯度；使用优化方法更新参数。随后重复这一过程，直至达到停止条件。与第3章的梯度下降相比，网络的参数更多，梯度也需要沿网络结构逐层计算。

如果训练损失下降而验证指标不再改善，可能出现过拟合；如果两者都很差，可能是模型表示、优化配置或特征处理不足。MLP 尤其需要合适的输入尺度、随机初始化控制和合理的训练轮数。

\subsection{随机初始化为什么不能全部取一样的值}
若同一层所有隐藏单元具有相同输入权重、相同偏置，并在后续连接中保持对称，它们会得到相同的前向结果和对称的梯度，更新后仍相同。名义上有许多单元，实际可能重复同一个表示。随机初始化的一个作用是打破这种对称，使不同单元有机会承担不同功能。

随机不等于越大越好。过大的初始权重可能把Sigmoid推入饱和区，或让激活和梯度尺度异常；过小也可能使信号很弱。成熟软件使用与层大小、激活有关的初始化策略。读者首先要固定随机种子以复现，再改变若干预先规定的种子检查结论，不挑选最好种子作为唯一报告。

\subsection{损失、惩罚与评价指标的对应}
对回归，若采用平均半平方误差加 $\lambda\|W\|_2^2/2$，训练包含预测误差与参数惩罚。偏置是否惩罚、对哪些层求和以及按多少样本归一化，依赖明确目标。sklearn的 \texttt{alpha} 与本书自定义 $\lambda$不能在没有核对目标尺度时机械等同；本书小网络手算使用写明的目标，不把手算学习率直接当成全部软件默认。

对分类，用真实类别对应概率的负对数计算交叉熵。准确率只统计最大概率类别是否正确，交叉熵还区分概率支持程度，因此两者可能呈现不同变化。训练准确率达到1以后，训练损失仍可能下降；这并不证明对新样本的概率越来越可靠。

\subsection{训练预算与停止原因}
一次参数更新处理一个批次；一个epoch通常遍历训练样本一次；\texttt{max\_iter}在不同求解器中的含义应依据接口说明。SGD或Adam通常按轮次组织，L-BFGS的迭代不是同一种逐批更新。因此，比较“都训练200次”时要说明比较的是何种计数。

收敛警告表示算法未达到声明的停止准则，不意味着所有预测都无效，也不允许直接忽略。应记录警告，检查输入尺度、目标、损失变化和预算，决定增加预算或调整配置。只有在误差和停止条件同时得到理解时，才有理由把模型作为完成的训练结果。

\section{反向传播与参数更新}
\textbf{反向传播（backpropagation）}的作用是高效计算损失对多层参数的梯度。它使用链式法则，从输出层的误差影响出发，逐层计算前面各层参数的导数。\textbf{优化器}则依据这些梯度及学习率等配置更新参数。前者负责“梯度怎么算”，后者负责“参数怎样变”，二者不能混为一谈。

考虑一个只有一个输入、一个隐藏单元、一个输出的简化网络，令
\begin{equation}
h=\phi(w_1x+b_1),\qquad \hat y=w_2h+b_2,
\end{equation}
并用 $\ell=\frac12(\hat y-y)^2$。通过链式法则，
\begin{equation}
\frac{\partial\ell}{\partial w_2}=(\hat y-y)h,
\qquad
\frac{\partial\ell}{\partial w_1}=(\hat y-y)w_2\phi'(w_1x+b_1)x.
\end{equation}
这里先求输出与目标的差，再把影响经输出权重和激活函数导数传回隐藏层。这就是反向传播在小网络上的基本机制。上述推导是为了理解计算路径，而不是要求学生手写完整深层网络训练器。

\subsection*{一次有具体数值的反向传播}
取单样本$x=2$、目标$y=1$，当前参数$w_1=0.5,b_1=0,w_2=3,b_2=0$，隐藏层仍用ReLU。前向得到$h=1$、$\hat y=3$、残差$e=2$、损失$\ell=2$。隐藏单元的输入为正，ReLU在此处的导数为1。由链式法则，四个梯度依次为
\[
\frac{\partial\ell}{\partial w_2}=2\times1=2,\quad
\frac{\partial\ell}{\partial b_2}=2,\quad
\frac{\partial\ell}{\partial w_1}=2\times3\times1\times2=12,\quad
\frac{\partial\ell}{\partial b_1}=2\times3\times1=6.
\]
注意损失对隐藏参数的影响要经过输出权重$w_2$；只用$e\times x$会漏掉这条计算路径。

用学习率0.01，\textbf{所有梯度都按同一组旧参数计算}，同时更新得到$w_1=0.38,b_1=-0.06,w_2=2.98,b_2=-0.02$。重新前向计算，$h=0.70$、$\hat y=2.066$，新损失为$\tfrac12(2.066-1)^2=0.568178$，小于原来的2。这个例子检验了当前一步，并不保证任意学习率下都下降。ReLU在零点不可微，本例特意避开零点，实际实现使用规定的处理方式。

\begin{lstlisting}[caption={核对一次网络参数更新}]
x, y = 2., 1.
w1, b1, w2, b2 = 0.5, 0., 3., 0.
h = max(0., w1*x+b1)
error = w2*h+b2-y
active = float(w1*x+b1 > 0.)
grad = (error*w2*active*x, error*w2*active, error*h, error)
w1, b1, w2, b2 = [p-0.01*g for p,g in
                  zip((w1,b1,w2,b2), grad)]
new_pred = w2*max(0., w1*x+b1)+b2
print(new_pred, 0.5*(new_pred-y)**2)
\end{lstlisting}


梯度得到以后，最简单的更新仍是 $\theta\leftarrow\theta-\eta\nabla_\theta\ell$；实际模型也可使用带动量的方法或 Adam 等优化器。

动量方法在当前梯度之外参考过去更新的累积方向；Adam则利用梯度历史的统计量调整各参数的更新尺度。它们仍然需要由反向传播提供梯度，不是替代了梯度计算。这里先理解职责；不要求把所有优化器推导都放进同一次网络入门实验。
批量大小、学习率和正则化会共同影响收敛。观察损失曲线时应分清横轴是更新次数、训练轮数还是已处理样本数。

\subsection{从局部导数到整条计算路径}
设某参数影响一个隐藏量，隐藏量又影响多个输出，那么损失对该参数的导数要累加所有路径的贡献。反向传播把计算图中已经得到的局部导数重复利用，从输出向输入方向传递。它是链式法则的高效组织，不是一种另外规定“正确方向”的经验规则。

对单隐藏层、Softmax与交叉熵的组合，输出得分梯度简化为 $\boldsymbol\delta_2=\mathbf p-\mathbf y$，这里 $\mathbf y$是目标的独热向量。于是
\[
\frac{\partial L}{\partial W_2}=\boldsymbol\delta_2\mathbf h^\mathsf T,
\qquad
\boldsymbol\delta_1=(W_2^\mathsf T\boldsymbol\delta_2)\odot\phi'(\mathbf z),
\qquad
\frac{\partial L}{\partial W_1}=\boldsymbol\delta_1\mathbf x^\mathsf T.
\]
$\odot$表示对应元素相乘。各偏置梯度分别为对应的 $\boldsymbol\delta$。一批样本时累加再按目标定义取平均；有正则化时还需加上其导数。逐项检查形状，就能判断某次矩阵乘法是否把样本维度与单元维度混淆。

这些式子是对本章单样本手算的概括，不要求读者从零实现完整深度框架。关键是理解输出误差怎样经过权重与激活导数影响隐藏参数，以及为什么所有梯度应基于同一次前向计算对应的旧参数。

\subsection{用数值差分检查反向计算}
对于可微点，选小的 $\epsilon$，比较解析梯度与
\[
\frac{L(\theta_j+\epsilon)-L(\theta_j-\epsilon)}{2\epsilon}.
\]
本章单隐藏单元例子的梯度为 $(12,6,2,2)$，分别对应 $w_1,b_1,w_2,b_2$。配套程序逐个扰动参数，数值差分与这些值一致，再用学习率0.01更新，核对损失从2降到0.568178。这一检查独立于链式计算代码，因此能帮助发现漏乘权重、符号错误或梯度使用了新参数的问题。

差分太大时近似不够局部，太小时受浮点消去影响；在ReLU折点附近，左右导数可能不同，也不能机械要求中心差分等于某个约定导数。本例故意使用正预激活，避免把不可微点处理问题与基础链式计算混在一起。

\subsection{计算梯度与选择更新方法是两件事}
普通梯度下降直接使用当前梯度；带动量的方法结合过去方向；Adam还利用梯度历史的尺度信息；L-BFGS利用拟牛顿近似选择更新。它们可以共享相同损失与反向传播梯度，但路径和停止标准不同。不能说“反向传播就是梯度下降”，也不能说“神经网络只能用普通梯度下降训练”。

对小型表格数据，L-BFGS是值得比较的选择；对较多记录，按小批次更新可能更合适。选择需要结合数据规模、目标曲面与资源，而不是把优化器名称作为新模型类别。不同算法得到不同局部解时，还需通过验证判断，而不是只按最后训练损失选择。

\section{使用 MLPClassifier 完成分类}
Wine 有三个类别和13个数值特征，前面已经用于逻辑回归、近邻与 SVM。此处继续使用它，模型复杂性是主要变化。由于 MLP 对输入尺度敏感，必须把 \texttt{StandardScaler} 放入 \texttt{Pipeline}；否则具有较大数值尺度的特征可能不利于优化。

\begin{lstlisting}[caption={Wine：逻辑回归与 MLP 的统一交叉验证比较}]
from sklearn.datasets import load_wine
from sklearn.model_selection import StratifiedKFold, cross_val_score
from sklearn.preprocessing import StandardScaler
from sklearn.pipeline import make_pipeline
from sklearn.linear_model import LogisticRegression
from sklearn.neural_network import MLPClassifier

X, y = load_wine(return_X_y=True)
from sklearn.model_selection import train_test_split
X, X_hold, y, y_hold = train_test_split(
    X, y, test_size=0.2, stratify=y, random_state=42
)
# All comparisons below use development X, y only.
cv = StratifiedKFold(5, shuffle=True, random_state=7)
models = {
    "logistic": make_pipeline(
        StandardScaler(), LogisticRegression(max_iter=3000)),
    "mlp": make_pipeline(
        StandardScaler(), MLPClassifier(hidden_layer_sizes=(20,),
                                       alpha=0.01, max_iter=800,
                                       random_state=42)),
}
for name, model in models.items():
    scores = cross_val_score(model, X, y, cv=cv, scoring="accuracy")
    print(name, scores.mean(), scores.std())
\end{lstlisting}

这段代码给出公平比较的基础形式。若要把它作为正式课程模型选择，需要在冻结的开发数据上运行，并保留最终测试集。如果某次迭代尚未满足停止条件，sklearn 可能给出收敛警告；这种警告需要检查训练曲线与参数配置，而不应被静默忽略。模型训练时的随机初始化、样本顺序和优化器细节都可能影响结果。

\fig{ch12_mlp_loss.pdf}{Wine 上使用固定随机状态训练的小型 MLP 的训练损失曲线。该曲线属于训练目标，不代表同一横轴上独立验证准确率的变化。}

\subsection{一套可复现的宽度与初始化比较}
在106条葡萄酒训练记录上拟合标准化和单隐藏层网络，在36条开发验证记录上评价。预先比较隐藏宽度4、12、32，每种使用种子0到4，固定L-BFGS与 \texttt{alpha=0.1}。15次运行的训练准确率都为1，验证准确率都为 $35/36\approx0.9722$。这组结果没有显示准确率随宽度或种子变化，不能为了表现“网络随机性”而绘制不存在的准确率波动。

但各次验证交叉熵与求解迭代数确实不同。宽度4的验证交叉熵约在0.0573至0.0712之间，宽度32约在0.0544至0.0668之间；例如宽度32、种子3的迭代数为140，而同宽度种子2为34。相同类别预测不意味着相同概率和优化路径。

这种小样本结果提示两点：准确率只能以 $1/36$ 为步进变化，区分细微改善的能力有限；不同随机状态可以在指标平台上仍产生不同内部模型。报告应保留所有预先规定的运行，而非仅选最低交叉熵或最短运行的一次。

\subsection{网络与逻辑回归比较的合理顺序}
先确认逻辑回归基线是否已很好地区分类别，再问隐藏层是否带来稳定收益。输入尺度、训练划分和目标应一致；网络可以需要不同的优化预算，但不能偷偷增加基线没有的特征。若网络没有明显改善，也是一项有效结果：当前任务可能不需要额外表示能力，或当前样本不足以支持复杂模型。

代码运行时，\texttt{hidden\_layer\_sizes=(8,)}表示一个8单元隐藏层，\texttt{(8,4)}表示两个隐藏层而不是12个并列单元。\texttt{activation}主要控制隐藏层非线性，输出层形式由分类或回归任务决定。\texttt{predict\_proba}返回类别顺序必须与 \texttt{classes\_}核对，不能把第2列未经检查就称为某个科学类别的概率。

\section{网络规模、正则化与训练诊断}
隐藏层规模可以写为 \texttt{hidden\_layer\_sizes=(20,)}，表示一个有20个单元的隐藏层；\texttt{(20, 10)} 则表示两层，分别有20和10个单元。层数和单元数增加通常提高模型表达能力，也增加优化与过拟合控制的负担。不能简单以“层数越多越好”判断配置。

课堂实验建议先固定其他配置，仅改变隐藏层规模；再在固定规模下改变正则化参数 \texttt{alpha} 或最大训练轮数。用于观察的损失曲线一般来自已拟合模型的 \texttt{loss\_curve\_} 属性；不同求解器对应的属性与训练行为需要查阅 API。若启用早停，内部会使用部分训练样本作验证，必须说明其与外部交叉验证和最终测试的关系。

\begin{lstlisting}[caption={观察 MLP 训练损失与初始化稳定性}]
from sklearn.pipeline import make_pipeline
from sklearn.preprocessing import StandardScaler
from sklearn.neural_network import MLPClassifier
from sklearn.datasets import load_wine

X, y = load_wine(return_X_y=True)
from sklearn.model_selection import train_test_split
X, X_hold, y, y_hold = train_test_split(
    X, y, test_size=0.2, stratify=y, random_state=42
)
# All comparisons below use development X, y only.
for seed in [1, 7, 42]:
    pipe = make_pipeline(
        StandardScaler(),
        MLPClassifier(hidden_layer_sizes=(16,), solver="adam",
                      alpha=0.01, max_iter=500, random_state=seed)
    )
    pipe.fit(X, y)
    net = pipe.named_steps["mlpclassifier"]
    print(seed, len(net.loss_curve_), net.loss_curve_[-1])
\end{lstlisting}

这里使用前面划定的开发部分训练，只是为了观察优化轨迹，不应把它产生的训练损失视为泛化评价；对照实验的预测性能仍应通过开发/验证方案确定。重复随机种子也不能替代新样本评价，它只检查某部分训练随机性是否明显影响计算。

\subsection*{从实验中辨认不同问题}
若训练损失持续下降而验证分数恶化，要考虑过拟合；若训练曲线没有明显下降，要检查学习率、尺度、激活函数和数据问题；若训练分数很高但特定类别召回率低，要回到样本分布、目标定义和评价指标。模型复杂度与性能之间没有单一方向的必然联系。

\subsection{用训练轨迹区分问题，而不是一次调整所有参数}
本轮另设一个12单元网络，以标准化后的训练记录执行固定轮数的SGD更新，每轮记录训练交叉熵和外部验证交叉熵。标准化只在训练记录上拟合，外部验证不用于梯度计算。图中每个点代表完成一轮后的检查，不能与每个小批次的瞬时损失混用。
\fig{v3_ch12_training.pdf}{葡萄酒小型网络的训练与外部验证交叉熵轨迹。该图检验优化与泛化行为，固定训练预算并完整保留轨迹。}

若训练与验证损失都高且仍明显下降，首先考虑预算是否不足；若训练高且几乎不变，检查尺度、学习率、梯度与表示；若训练继续下降而验证明显变差，考虑复杂度、正则化和停止位置。若两者都已较低但某类召回率仍差，要检查总体指标是否掩盖了少数类，而不是只增加网络宽度。

诊断顺序应一次只改变少量因素。首先验证数据和目标正确，再比较缩放，接着检查优化预算，最后研究网络结构和正则化。若同时改宽度、学习率、批大小和训练数据，分数变化很难归因，学生也难以形成可迁移的经验。

\subsection{早停、内部验证与重复对象}
软件的 \texttt{early\_stopping} 可能从输入训练数据中另划一部分作内部验证。应弄清它监测的是准确率、$R^2$还是损失，以及多久不改善才停止。对包含同一对象多次观测的数据，内部随机拆行可能与外部对象分组要求不一致。此时可通过外部受控验证和固定训练轮次选择复杂度，避免让隐藏的切分破坏研究设计。

在Pipeline中，外层标准化可能先在传给估计器的全部训练数据上拟合，再由网络内部划分早停集。这不涉及外部最终测试，但内部早停评价与严格“只在内部训练子集学习缩放”的流程并不完全相同。科学任务要求严格内部独立时，应把这些步骤显式组织起来，而不是认为用了Pipeline就解决了所有层级的信息问题。

\subsection{参数量不等于有效数据量}
139个参数面对106条训练记录，并不自动意味着模型不可用，也不意味着有足够信息估计每个参数的物理含义。正则化、结构、输入相关性和任务简单程度共同影响泛化。更宽的网络可以存在许多实现相似函数的参数配置；因此不应把单个权重的数值大小直接当作物理重要性。

\section{MLPRegressor 与后续神经网络学习}
MLP 也可用于回归。与分类网络相比，回归输出层一般用于给出连续数值，不需要把单个输出压到类别概率区间。教学中可以使用 \texttt{MLPRegressor} 对 California Housing 进行小规模回归演示，比较线性回归和森林的验证误差，但应控制训练轮数与数据规模，避免把课程时间主要花在调参上。

\begin{lstlisting}[caption={回归 MLP 的基本流程，需使用第5章的住房数据}]
from sklearn.neural_network import MLPRegressor
from sklearn.preprocessing import StandardScaler
from sklearn.pipeline import make_pipeline

reg = make_pipeline(
    StandardScaler(),
    MLPRegressor(hidden_layer_sizes=(32, 16),
                 alpha=0.01, max_iter=500, random_state=42)
)
# Use development cross-validation before final evaluation.
\end{lstlisting}

\subsection*{提高练习的数据：手写数字的小图像}
\label{sec:digits-intro}
\textbf{手写数字数据集（Digits）}包含1797幅$8\times8$灰度图像，目标是图像表示的0至9中的一个数字类别。\textbf{像素}是图像网格中的一个单元，这组数据每格的数值为0至16，表示整理后的灰度强度。每幅图像共有64个像素，因此可用64个数表示一条样本。\textsuperscript{[28]}

按第一行、第二行依次排开8行像素，就得到长度64的向量，这个操作称为\textbf{展平}。只要保存原来的行列顺序，就能再把它恢复成$8\times8$网格；展平并不丢弃像素数值。标签0至9是图像内容的类别，虽然写成数字，也不是要回归计算“数字8比数字4大多少”。

\begin{lstlisting}[caption={查看图像数组与展平后的同一条样本}]
from sklearn.datasets import load_digits
import numpy as np
digits = load_digits()
print(digits.images.shape, digits.data.shape)
print(digits.images[0])
print(digits.target[0])
assert np.array_equal(digits.images[0].ravel(), digits.data[0])
\end{lstlisting}
两个形状分别为$(1797,8,8)$和$(1797,64)$。\texttt{images[0]}是第一幅图的网格，\texttt{data[0]}是同一幅图的向量，\texttt{target[0]}是已知数字类别。先在纸上追踪前两行像素怎样变成向量前16项，再训练分类器，能够避免只看一串64维数字而不知道输入来自哪里。

\fig{ch12_digits.pdf}{sklearn Digits 中若干 $8\times8$ 手写数字图像。每幅图具有64个像素特征，可展平为网络输入，用于提高练习。}

Digits 数据可以进一步展示图像如何变为数值特征，展平仍保留全部像素值及固定位置顺序，但普通全连接 MLP 没有显式施加局部邻域连接与共享权重的结构。后续卷积网络会用更专门的结构处理空间邻域，循环或注意力结构可用于其他类型输入。本书只作学习方向概述，不将大型网络训练、复杂架构或当前生成式 AI 作为本章必做实验。

\subsection*{专业迁移：分类模型是否值得更复杂？}
沿用第6章的 SDSS 或 Sonar 开发数据，对照逻辑回归、一种树集成模型与 MLP。记录同一数据条件下的分类表现、训练稳定性、主要错误类别和资源需求。若 MLP 未取得改善，可能是训练配置、数据规模、特征或任务本身的原因，不能仅凭一次实验否定神经网络。若取得改善，也不能忽略模型复杂度与科学解释的额外成本。

\scientificnote{监督训练时使用的标签可能带选择效应、观测误差或类别边界不确定性。网络能较好拟合标签分布，并不等于已经验证了这些类别的天体物理机制。}

\subsection{回归输出、目标尺度与误差单位}
MLP回归通常使用数值输出和相应的回归损失。若目标从几十分之一变为几万，梯度与正则化的相对尺度也会改变。可以在训练数据上学习目标缩放，预测后再逆变换；scikit-learn的 \texttt{TransformedTargetRegressor}可把这一步与模型组织在一起。目标缩放参数不能利用最终测试目标。

例如训练目标均值100、标准差20，网络在标准化目标上预测0.5，原单位预测为110。若标准化RMSE为0.2，则原单位为4，但只有同一线性缩放和对应评价记录下才能这样换算。对数变换不具有这种简单比例关系，逆变换后的均值偏差和评价需另行检查。

输入尺度与目标尺度作用不同：前者影响各输入方向和优化条件，后者影响输出误差与梯度。两者都不应在每一条新样本上重新拟合。训练后的模型应保存变换参数、特征顺序和全部网络权重，才能真正复现预测。

\subsection{图像展平之后，哪些信息保留、哪些结构没有利用}
本章提高练习使用已经介绍的8乘8手写数字图像。展平后64个像素仍保存各位置的数值，但普通全连接MLP并没有强制相邻像素共享处理规则。若打乱像素列并对训练和预测一致使用相同排列，全连接模型理论上可以重新学习相应权重；它不像卷积网络那样显式利用局部邻接和权重共享。

这解释了后续学习更复杂网络的动机：不是因为MLP不属于神经网络，而是因为不同数据结构可以支持不同的归纳偏好。对光谱序列、时序或图像，先理解数据轴的物理含义，再决定是否需要专门结构。不能把任意二维表格转成一张图后就声称获得了图像物理邻接。

\subsection{进入后续课程前应掌握的最小闭环}
学生应能画出一个小网络，标明张量形状，完成一次前向计算、损失计算、梯度计算和更新；能区别表示不足、优化困难与过拟合；能在同一专业任务上与线性及树模型公平比较。具备这些能力后，再进入卷积、注意力或大型预训练模型，概念才有坚实基础。

这章不要求学生把每个高阶优化器都推导一遍，也不以网络参数量作为学习成果。用一个可解释的小模型完成可靠实验，比运行一个无法解释输入、标签和评价的大网络更符合本书的科学建模目标。

\chapterpractice{必做：Wine 的 StandardScaler + MLPClassifier，与逻辑回归在固定开发/验证方案下比较；改变隐藏层规模、正则化和随机种子，并画训练损失曲线。提高：Digits 分类和误分类图；回归：使用住房数据尝试 MLPRegressor；专业迁移：选 SDSS 或 Sonar 比较三种有代表性的监督模型。}
\chapterexercise
\begin{enumerate}
\item 从逻辑回归表达式出发，说明加入一个隐藏层后模型结构有哪些变化。
\item 如果多层网络中完全不使用非线性激活，多个线性层的组合具有何种表达能力？
\item 对一个单隐藏单元网络，写出前向计算步骤，并解释链式法则怎样把损失影响传回隐藏层。
\item 反向传播与梯度下降分别负责什么？\texttt{fit()} 是否保证每种 sklearn 模型都使用这两种方法？
\item 仅依据 \texttt{loss\_curve\_} 持续下降，能否证明 MLP 在未见数据上的分类效果改善？为什么？
\item 请设计一个控制变量实验，比较 Wine 上的两种隐藏层规模。你需要固定什么？报告哪些指标？
\item 在 Sonar 的小样本场景中，哪些因素可能使不同 MLP 随机初始化结果波动较大？
\end{enumerate}

\section*{深化练习与方法检验}
本组题强调解释、计算与验证；D编号在本章内独立计数。计算题应写出所用定义和中间步骤，实验题保留实际输出，包括不支持预期的结果。
\begin{enumerate}[label=D\arabic*.]
\item 一个具有13个输入、8个隐藏单元和3个输出单元的全连接网络，包含各层偏置，共有多少可学习标量？若再插入4单元隐藏层，参数数目如何变化？
\item 用 $h_1=\max(0,x_1-x_2)$、$h_2=\max(0,x_2-x_1)$ 和输出 $h_1+h_2$ 计算四个二值输入。解释它怎样表示XOR，同时说明这个例子没有证明什么。
\item 对正文单隐藏单元反向传播例子，用中央差分近似四个梯度。为什么步长不宜无限小，ReLU输入为零附近又会带来什么问题？
\item 同一Wine划分上多个网络宽度和种子都得到35/36的验证准确率，但交叉熵不同。应该报告什么？能否从一次结果确定较宽网络具有显著优势？
\end{enumerate}
